% ----------------------
\subsection{Section 50 (9 May 1831)}
% ----------------------
	\label{DC50}
	
	% -----
	\subsubsection{Early versions of this section}
	% -----
	
		Early versions of this section appear in both \hyperref[EXSS-DC50-EARLY-GILBERT]{Sidney Gilbert's notebook} and the \hyperref[EXSS-DC50-EARLY-NOTEBOOK]{Hyde and Smith notebook}, in addition to the JSP. I've drawn from these versions in some of my notes below.
	
	% -----
	\subsubsection{Historical Background}
	% -----
		\label{DC50-HB}
		
		% --
		\paragraph{JSP}
		% --
		
			``In late fall 1830, Oliver Cowdery and his missionary companions left Kirtland, Ohio, on their way to preach to the American Indians in the territory just west of Missouri after having baptized more than one hundred persons into the church. In addition, several leaders among the Ohio converts departed as well. Frederick G. Williams, who accompanied Cowdery on the mission to the Indians, departed in November, and Edward Partridge and Sidney Rigdon traveled to New York soon afterward to meet with JS. The new church members in Ohio were left without an experienced leader until John Whitmer arrived in January 1831. Upon his arrival, Whitmer noted with dismay that `the enemy of all righteous had got hold of some of those who profesed to be his followers, because they had not sufficent knowledge to detect him in all his devices.' When Parley P. Pratt, one of the missionaries who introduced the teachings of the church to Ohio residents, returned to Kirtland from Missouri in spring 1831, he too was concerned with `new and strange' behaviors he saw in the church: `As I went forth among the different branches, some very strange spiritual operations were manifested, which were disgusting, rather than edifying. Some persons would seem to swoon away, and make unseemly gestures, and be drawn or disfigured in their countenances. Others would fall into ecstacies, and be drawn into contortions, cramp, fits, etc. Others would seem to have visions and revelations, which were not edifying, and which were not congenial to the doctrine and spirit of the gospel. In short, a false and lying spirit seemed to be creeping into the Church.'
			
			``Such ostentatious displays did not go unnoticed in the surrounding community. The \textit{Painesville Telegraph} reported, `Immediately after Mr. R[igdon] and the four pretended prophets left Kirtland, a scene of the wildest enthusiasm was exhibited, chiefly, however, among the young people; they would fall, as without strength, roll upon the floor, and, so mad were they that even the females were seen in a cold winter day, lying under the bare canopy of heaven, with no couch or pillow but the fleecy snow.'
			
			``After JS arrived in Ohio in February 1831, he worked to curtail what he perceived to be excessive and ostentatious spiritual behaviors among the believers. In early March he wrote to his brother Hyrum that he had been `ingageed in regulating the Churches here as the deciples are numerous and the devil had made many attempts to over throw them it has been a Serious job but the Lord is with us and we have overcome and have all things regular.' A revelation a few days later urged church members to walk `uprightly before me...that ye may not be seduced by evil spirits or doctrines of Devils or the commandments of men for some are of men \& others of Devils Wherefore beware lest ye are deceived.'
			
			``The problem of discerning between appropriate and inappropriate behaviors apparently continued for several months. When John Whitmer described the 9 May revelation featured here in his later history, he provided this setting for it:
			
			\begin{quote}
				Some had visions and could not tell what they saw, Some would fancy to themselves that they had the sword of Laban, and would wield it as expert as a light dragoon, some would act like an Indian in the act of scalping, some would slide or scoot and [on] the floor, with the rapidity of a serpent, which the[y] termed sailing in the boat to the Lamanites, preaching the gospel. And many other vain and foolish manoeuvers that are unseeming, and unprofitable to mention. Thus the devil blinded the eyes of some good and honest disciples...These things grievd the servants of the Lord, and some conversed together on this subject, and otheers came in and we were at Joseph Smith Jr. the seers, and made it a matter of consultation, for many would not turn from their folly, unless God would give a revelation, therfore the Lord spake to Joseph.
			\end{quote}
			
			``Surviving accounts demonstrate that in the months following this revelation, the elders of the church employed the revelation's guidelines instructing them what to do `if ye behold a spirit manifested that ye cannot understand.' JS was involved in one such situation a few weeks after this revelation, when he determined that a manifestation involving Harvey Whitlock was of the devil.%
				\footnote{%
					% \label{}%
					JSP note: ``According to Levi Hancock, during a conference held in early June 1831, `Joseph put his hands upon Harvey Whitlock and ordained him to the high Priesthood he [Whitlock] turned as black as Lymon [Lyman Wight] was white his fingers was set like Claws he went round the Room and showed his hands and tryed to speak his eyes was in the shape of Ovil Oes [oval \textit{o}'s] Hyram smith said Joseph that is not God Joseph said do not speak against this[.] I will not beleive said Hyrum unless you inquire of God and owns it Joseph bowed his head a short time and got up and commanded satan to leave Harvey laying his hands upon his head at the same time.' (Hancock, Autobiography, 90.) '' The info for the Autobiography is ``Hancock, Levi. Autobiography, ca. 1854. Photocopy. CHL. MS 8174.''
					}
			A short time later, Jared Carter recounted another attempt to follow the specific directions given in the revelation when he encountered a woman possessed with a spirit during a church meeting.''%
				\footnote{%
					% \label{}%
					JSP note: ``Carter wrote of an experience in Amherst, Ohio: `At length I proved by a revelation that had ben given to the Elders concerning spiriths that these spirits visionary exercises as they were called were not of the Lord...as we was about to atend to the administration of the communion there was a young woman taken with an exercise that brought her on to the floor \& be cause I doubted of such maner of influences in a public congragation, I reqested Brother Silvester [Sylvester Smith] that we should try that spirit acording to the revelation that god had given he immediately complied with my request we then neeled down and asked our heavenly father in the name of Christ that if that spirit that that sister possesed was of him that he would give it to us \& we prayed in faith but we did not receive that Spirit we then arose \& I sat apon my Seat Silent for some minutes but Brother Sylvester arose and laid hands apon the Sister but this was not as the commandment directs.' Carter viewed the revelation featured here as a formula for casting out evil spirits, and his journal continues: `The command reads thus [``]wherefore it shal come to pass that if you behold a spirit manifest that ye cannot understand \& ye receive not that Spirit ye Shall ask the Father in the name of Jesus \& if he give you not that Spirit then ye may know that it is not of god \& it Shall be given un to you power over that Spirit \& ye shall proclaim against that Spirit with aloud voice that it is not of god not with railing accusation that ye be not overcome neither with bosting nor rejoising lest you be Seized there with[''] now after Silvester had made some communication which was not propclameing against the Spirit as I believeed it had ought to have been that is against the spirit that we had prayed concerning I then arose \& proclaimed against that spirit with a loud voice...and from that time forward that spirit never came in to the meeting when I was preasant in this display of the power of God I had one of the most infalable proofs of the divine origen of the above mentioned revalation.' (Carter, Journal, 24--26, 29.)'' The info for the journal is: ``Carter, Jared. Journal, 1831--1833. CHL. MS 1441.'' See HDDC 645--647 for a lengthier excerpt from Carter's journal discussing these issues.
					}
		
	% -----
	\subsubsection{Versions}
	% -----
		\label{DC50-V}
		
		See the \href{https://drive.google.com/file/d/1goAvNz8jS4R8slYfMG_db55Ty2z_vmgK/view?usp=sharing}{sections comparison} document.
	
		\begin{compactdesc}
			\item[JSP] \hyperref[EMS-1831-JS-9-MAY-1831]{9 May 1831}
			\item[EMS] 1:3 (August 1832), 17
			\item[BC] Chapter 53, 119--123
			\item[DC 1835] Section 17, 134-136
			\item[TS] 5:2 (15 January 1844), 402--403
			\item[MS] 5:6 (November 1844), 84--85
			\item[DC 1844] Section 65, 190--194
		\end{compactdesc}

	% -----
	\subsubsection{Section 50:1} % *DC50-1 
	% -----
		\label{DC50-1}

		\hyperref[HEARKEN]{HEARKEN}, O ye elders of my church, and give ear to the voice of the living God; and attend to the words of wisdom%
			\footnote{%
				% \label{}%
				See the phrase ``word(s) of wisdom'' \hyperref[WISDOM-PHRASES]{here}.
				}
		which shall be given unto you, according as ye have asked and are agreed as touching the church, and the spirits which have gone abroad in the earth.

		% \begin{framed}
		% \scriptsize
			% \begin{compactdesc}
				% \item[JSP] 
				% % \item[BC] 
				% % \item[]
			% \end{compactdesc}
		% \end{framed}

	% -----
	\subsubsection{Section 50:2} % *DC50-2 
	% -----
		\label{DC50-2}

		Behold, verily I say unto you, that there are many spirits which are false spirits, which have gone forth in the earth, deceiving the world.%
			\footnote{%
				% \label{}%
				Compare \hyperref[DC46-7]{DC 46:7}.
				}

		% \begin{framed}
		% \scriptsize
			% \begin{compactdesc}
				% \item[JSP] 
				% % \item[BC] 
				% % \item[]
			% \end{compactdesc}
		% \end{framed}

	% -----
	\subsubsection{Section 50:3} % *DC50-3 
	% -----
		\label{DC50-3}

		And also Satan hath sought to deceive you, that he might overthrow you.

		% \begin{framed}
		% \scriptsize
			% \begin{compactdesc}
				% \item[JSP] 
				% % \item[BC] 
				% % \item[]
			% \end{compactdesc}
		% \end{framed}

	% -----
	\subsubsection{Section 50:4} % *DC50-4 
	% -----
		\label{DC50-4}

		Behold, I, the Lord, have looked upon you, and have seen abominations in the church that profess my name.

		% \begin{framed}
		% \scriptsize
			% \begin{compactdesc}
				% \item[JSP] 
				% % \item[BC] 
				% % \item[]
			% \end{compactdesc}
		% \end{framed}

	% -----
	\subsubsection{Section 50:5} % *DC50-5 
	% -----
		\label{DC50-5}

		But blessed are they who are faithful and endure, whether in life or in death, for they shall inherit eternal life.

		% \begin{framed}
		% \scriptsize
			% \begin{compactdesc}
				% \item[JSP] 
				% % \item[BC] 
				% % \item[]
			% \end{compactdesc}
		% \end{framed}

	% -----
	\subsubsection{Section 50:6} % *DC50-6 
	% -----
		\label{DC50-6}

		But wo unto them that are deceivers and hypocrites, for, thus saith the Lord, I will bring them to judgment.

		% \begin{framed}
		% \scriptsize
			% \begin{compactdesc}
				% \item[JSP] 
				% % \item[BC] 
				% % \item[]
			% \end{compactdesc}
		% \end{framed}

	% -----
	\subsubsection{Section 50:7} % *DC50-7 
	% -----
		\label{DC50-7}

		Behold, verily I say unto you, there are hypocrites among you, who have deceived some, which has given the adversary power;%
			\footnote{%
				% \label{}%
				See \hyperref[2NEPHI-2-29]{2 Nephi 2:29}.
				}
		but behold such shall be reclaimed;

		% \begin{framed}
		% \scriptsize
			% \begin{compactdesc}
				% \item[JSP] 
				% % \item[BC] 
				% % \item[]
			% \end{compactdesc}
		% \end{framed}

	% -----
	\subsubsection{Section 50:8} % *DC50-8 
	% -----
		\label{DC50-8}

		But the hypocrites shall be detected and shall be cut off, either in life or in death, even as I will; and wo unto them who are cut off from my church, for the same are overcome of the world.%
			\footnote{%
				% \label{}%
				See \hyperref[WORLD-PHRASES]{here} for instances of ``overcome'' and ``world'' appearing together---this is the only one that I could find which mentions the world overcoming \textit{us}, rather than the other way around. Compare especially \hyperref[DC50-35]{verse 35} and \hyperref[DC50-41]{verse 41} below. It also reminds me of \hyperref[DC76-29]{DC 76:29--30}---see those verses and notes there, as well as \hyperref[DC52-18]{DC 52:18}, 
				}

		% \begin{framed}
		% \scriptsize
			% \begin{compactdesc}
				% \item[JSP] 
				% % \item[BC] 
				% % \item[]
			% \end{compactdesc}
		% \end{framed}

	% -----
	\subsubsection{Section 50:9} % *DC50-9 
	% -----
		\label{DC50-9}

			\footnote{%
				% \label{}%
				Beginning with verse 9 here we get what I think is one of the most amazing passages in all of the DC, including one of the most important discussions of the epistemology of revelation. For more notes and connections on these verses than what is here, see my \hyperref[EPIST-REVELATION-DC]{epistemology of revelation} notes.
				}%
		Wherefore, let every man beware%
			\footnote{%
				% \label{}%
				Note that, of the three early manuscript versions, two have ``be aware'' here, including the JSP version (which is RB1).
				}
		lest he do that which is not in \hyperref[TRUTH]{truth} and \hyperref[RIGHTEOUSNESS]{righteousness}%
			\footnote{%
				% \label{}%
				Reminds me of \hyperref[MORONI-10-6]{Moroni 10:6} and \hyperref[GOODNESS-truth]{goodness/truth} stuff.
				}
		before me.

		\begin{framed}
		\scriptsize
			\begin{compactdesc}
				\item[JSP] let every man be aware lest he do that which is not in truth \& righteousness before me
				\item[Gilbert] Wherefore let every man beware lest he do that which is not in truth \& righteousness before me,
				\item[Hyde and Smith] wherefore let every man be aware lest he do that which is not in truth \& righteousness before me,
				\item[BC] 9 Wherefore, let every man beware lost [sic] he do that which is not in truth and righteousness before me.
			\end{compactdesc}
		\end{framed}

	% -----
	\subsubsection{Section 50:10} % *DC50-10 
	% -----
		\label{DC50-10}

		And now come, saith the Lord, by the Spirit, unto the elders of his church, and let us \hyperref[REASON]{reason} together, that ye may understand;%
			\footnote{%
				% \label{}%
				The purpose of reasoning together is for us to \textit{understand}.
				}

		\begin{framed}
		\scriptsize
			\begin{compactdesc}
				\item[JSP] \& now come saith the Lord by the spirit by the Elders of his Church \& let us reason together that ye may understand
				\item[Gilbert] and now Come saith the Lord, by the Spirit unto the Elders of his Church and let us reason together that ye may understand,
				\item[Hyde and Smith] And now come saith the Lord by the Spirit unto the Elders of his church \& let us reason together that ye may understand
				\item[BC] 10 And now come, saith the Lord, by the Spirit, unto the elders of his church, and let us reason together, that ye may understand:
			\end{compactdesc}
		\end{framed}

	% -----
	\subsubsection{Section 50:11} % *DC50-11 
	% -----
		\label{DC50-11}

		Let us \hyperref[REASON]{reason} even as a man \hyperref[REASON]{reasoneth} one with another%
			\footnote{%
				% \label{}%
				The Lord is speaking here, so the invitation he is making is to us to come and reason with him, ``even as a man reasoneth one with another.''
				}
		face to face.

		\begin{framed}
		\scriptsize
			\begin{compactdesc}
				\item[JSP] let us reason even as a man reasoneth one with another face to face
				\item[Gilbert] let us reason as a man reasoneth one with another, face to face---
				\item[Hyde and Smith] let us reason even as a man reasoneth one with another face to face,
				\item[BC] 11 Let us reason even as a man reasoneth one with another face to face:
			\end{compactdesc}
		\end{framed}

	% -----
	\subsubsection{Section 50:12} % *DC50-12 
	% -----
		\label{DC50-12}

		Now, when a man \hyperref[REASON]{reasoneth} he is understood of man, because he \hyperref[REASON]{reasoneth} as a man; even so will I, the Lord, \hyperref[REASON]{reason} with you that you may understand.%
			\footnote{%
				% \label{}%
				An extremely important verse---the Lord is telling us that he will reason with us \textit{as a man}, so that we can understand. Compare to other scriptures like \hyperref[2NEPHI-31-3]{2 Nephi 31:3} and \hyperref[DC1-24]{DC 1:24}.
				}

		\begin{framed}
		\scriptsize
			\begin{compactdesc}
				\item[JSP] now when a man reasoneth he under[stand]eth of man because he reasoneth as a man even so will I the Lord reason with you that you may understand
				\item[Gilbert] now when a man reasoneth he understandeth of man, because he reasoneth as a man even so will I the Lord reason with you, that you may understand,
				\item[Hyde and Smith] now when a man reasoneth he is understood of man because he reasoneth as a man even so will I the lord reason with you that you may understand
				\item[BC] 12 Now when a man reasoneth, he is understood of man, because he reasoneth as a man; even so will I the Lord reason with you that you may understand:
			\end{compactdesc}
		\end{framed}

	% -----
	\subsubsection{Section 50:13} % *DC50-13 
	% -----
		\label{DC50-13}

		Wherefore, I the Lord ask you this question---unto what were ye ordained?

		\begin{framed}
		\scriptsize
			\begin{compactdesc}
				\item[JSP] wherefore I the Lord asketh you this question unto what were ye ordained
				\item[Gilbert] wherefore I the Lord asketh you this question, unto what was ye ordained
				\item[Hyde and Smith] wherefore I the Lord asketh you this question unto what was ye ordained,
				\item[BC] 13 Wherefore I the Lord asketh you this question, unto what were ye ordained?
			\end{compactdesc}
		\end{framed}

	% -----
	\subsubsection{Section 50:14} % *DC50-14 
	% -----
		\label{DC50-14}

		To preach my gospel by the Spirit, even the \hyperref[COMFORTER]{Comforter} which was sent forth to teach the truth.

		\begin{framed}
		\scriptsize
			\begin{compactdesc}
				\item[JSP] to Preach my Gospel by the spirit even the comforter which was sent forth to teach the truth
				\item[Gilbert] to preach my Gospel by the Spirit, even the Comforter which was sent forth to teach the truth,
				\item[Hyde and Smith] to preach my Gospel by the spirit even the comforter which was sent forth to teach the truth
				\item[BC] 14 To preach my gospel by the Spirit, even the Comforter which was sent forth to teach the truth;
			\end{compactdesc}
		\end{framed}

	% -----
	\subsubsection{Section 50:15} % *DC50-15 
	% -----
		\label{DC50-15}

		And then received ye spirits which ye could not understand, and received them to be of God; and in this are ye justified?

		\begin{framed}
		\scriptsize
			\begin{compactdesc}
				\item[JSP] \& then received you spirits which ye could not understand \& received them to be of God \& in this are ye Justified?
				\item[Gilbert] \& then received ye the Spirits which ye could not understand, and received them to be of God--- and in this are ye justified?---
				\item[Hyde and Smith] \& then received ye spirits which ye could not understand \& rec\supersu{d}\supers{.} them to be of God \& in this are ye justified!
				\item[BC] and then received ye spirits which ye could not understand, and received them to be of God, and in this are ye justified?
			\end{compactdesc}
		\end{framed}

	% -----
	\subsubsection{Section 50:16} % *DC50-16 
	% -----
		\label{DC50-16}

		Behold ye shall answer this question yourselves; nevertheless, I will be merciful unto you; he that is weak among you hereafter shall be made strong.

		\begin{framed}
		\scriptsize
			\begin{compactdesc}
				\item[JSP] Behold ye shall answer this yourselves nevertheles I will be mercyfull unto you he that is weak among you hereafter shall be made strong
				\item[Gilbert] Behold ye shall answer this question of yourselves, nevertheless I will be merciful unto you, he that is weak among you hereafter shall be made strong---
				\item[Hyde and Smith] behold ye shall answer this question of yourselves nevertheless I will be merciful unto you he that is weak among you hereafter shall be made strong
				\item[BC] 15 Behold ye shall answer this question yourselves, nevertheless I will be merciful unto you: 16 He that is weak among you hereafter shall be made strong.
			\end{compactdesc}
		\end{framed}

	% -----
	\subsubsection{Section 50:17} % *DC50-17 
	% -----
		\label{DC50-17}

		Verily I say unto you, he that is ordained of me and sent forth to preach the word of truth by the Comforter, in the Spirit of truth,%
			\footnote{%
				% \label{}%
				Note that the JSP (RB1) and Hyde and Smith versions do not have ``in,'' making ``the Spirit of truth'' in apposition to ``the Comforter.'' This would put the verse in line with places like \hyperref[JOHN-14-17]{John 14:17} and \hyperref[JOHN-15-26]{John 15:26}. The occurrences in John are important because that is where you get the original discussion of \hyperref[PARAKLETOS]{\grl{παράκλητος}}. The phrase also appears in \hyperref[DC93]{DC 93}; see all the occurrences of ``spirit of truth'' \hyperref[SPIRIT-PHRASES]{here}.
				}
		doth he preach it by the Spirit of truth or some other way?

		\begin{framed}
		\scriptsize
			\begin{compactdesc}
				\item[JSP] verily I say unto you he that is ordained of me \& sent forth to preach the word of truth by the comforter the spirit of truth doth he preach it by the spirit of truth or some other way
				\item[Gilbert] verily I say unto you he that is ordained of me, \& sent forth to preach the word of truth by the Comforter in the Spirit of truth, doth he preach it by the Spirit of truth, or some other way,
				\item[Hyde and Smith] verily I say unto you he that is ordained of me \& sent forth to preach the word of truth by the comforter the spirit of truth doth he preach it by the spirit of truth or some other way
				\item[BC] 17 Verily I say unto you, he that is ordained of me and sent forth to preach the word of truth by the Comforter, in the spirit of truth, doth he preach it by the spirit of truth, or some other way? 
			\end{compactdesc}
		\end{framed}

	% -----
	\subsubsection{Section 50:18} % *DC50-18 
	% -----
		\label{DC50-18}

		And if it be by some other way it is not of God.

		\begin{framed}
		\scriptsize
			\begin{compactdesc}
				\item[JSP] \& if by some other way it be not of God
				\item[Gilbert] \& if by some other way, it be not of God,
				\item[Hyde and Smith] if by some other way it be not of God,
				\item[BC] and if by some other way, it be not of God.
			\end{compactdesc}
		\end{framed}

	% -----
	\subsubsection{Section 50:19} % *DC50-19 
	% -----
		\label{DC50-19}

		And again, he that receiveth the word of truth, doth he receive it by the Spirit of truth or some other way?

		\begin{framed}
		\scriptsize
			\begin{compactdesc}
				\item[JSP] \& again he that receiveth the word of truth doth he receive it by the spirit of truth or some other way
				\item[Gilbert] and again he that receiveth the word of truth, doth he receive it by the Spirit of truth or some other way,
				\item[Hyde and Smith] \& again, he that receiveth the word of truth doth he receive it by the spirit of truth or some other way
				\item[BC] 18 And again, he that receiveth the word of truth, doth he receive it by the spirit of truth, or some other way?
			\end{compactdesc}
		\end{framed}

	% -----
	\subsubsection{Section 50:20} % *DC50-20 
	% -----
		\label{DC50-20}

		If it be some other way it is not of God.

		\begin{framed}
		\scriptsize
			\begin{compactdesc}
				\item[JSP] if it be some other way it be not of God
				\item[Gilbert] if it be some other way it be not of God.---
				\item[Hyde and Smith] if it be some other way it be not of God.
				\item[BC] if it be some other way, it be not of God;
			\end{compactdesc}
		\end{framed}

	% -----
	\subsubsection{Section 50:21} % *DC50-21 
	% -----
		\label{DC50-21}

		Therefore, why is it that ye cannot understand and know, that he that receiveth the word by the Spirit of truth receiveth it as it is preached by the Spirit of truth?

		\begin{framed}
		\scriptsize
			\begin{compactdesc}
				\item[JSP] therefore why is it that ye cannot understand \& know that he <that> receiveth the word by the spirit of truth receiveth it as it is preached by the spirit of truth
				\item[Gilbert] Wherefore why is it that ye cannot understand \& know that he that receiveth the word by the Spirit of truth receiveth it as it is preached by the Spirit of truth,
				\item[Hyde and Smith] Therefore why is it that ye cannot understand \& know that he that receiveth the word by the spirit of truth receiveth it as it is preached by the spirit of truth
				\item[BC] 19 Therefore, why is it that ye can not understand and know that he that receiveth the word by the spirit of truth, receiveth it as it is preached by the spirit of truth?
			\end{compactdesc}
		\end{framed}

	% -----
	\subsubsection{Section 50:22} % *DC50-22 
	% -----
		\label{DC50-22}

		Wherefore, he that preacheth and he that receiveth,
		\footnote{%
			% \label{}%
			They are edified \textit{because they understand each other}.
			}%
		understand one another,%
			\footnote{%
				% \label{}%
				Compare \hyperref[GENESIS-11-7]{Genesis 11:7}, which is part of the Babel tower story. This could be a significant reference, when you think about ``language'' and reasoning together and coming to mutual understanding, along with the connections given in \hyperref[DC50-12]{verse 12}.
				}
		and both are \hyperref[EDIFICATION]{edified} and rejoice together.

		\begin{framed}
		\scriptsize
			\begin{compactdesc}
				\item[JSP] wherefore he that <preacheth \& he that> receiveth understandeth one another \& both are edified \& rejoice together
				\item[Gilbert] wherefore, he that preacheth \& he that receiveth understandeth one another and both are edified and rejoice together
				\item[Hyde and Smith] wherefore he that preacheth \& he that receiveth understandeth one another \& both are edifyed \& rejoice together
				\item[BC] 20 Wherefore, he that preacheth and he that receiveth, understandeth one another, and both are edified and rejoice together;
			\end{compactdesc}
		\end{framed}

	% -----
	\subsubsection{Section 50:23} % *DC50-23 
	% -----
		\label{DC50-23}

		And that which doth not \hyperref[EDIFICATION]{edify} is not of God, and is darkness.

		\begin{framed}
		\scriptsize
			\begin{compactdesc}
				\item[JSP] \& that which doth not edify is not of God \& is darkness
				\item[Gilbert] and that which doth not edify is not of God and is darkness,
				\item[Hyde and Smith] \& that which doth not edify is not of God \& is darkness,
				\item[BC] and that which doth not edify, is not of God, and is darkness:
			\end{compactdesc}
		\end{framed}

	% -----
	\subsubsection{Section 50:24} % *DC50-24 
	% -----
		\label{DC50-24}

		That which is of God is light;%
			\footnote{%
				% \label{}%
				Compare \hyperref[ALMA-32-35]{Alma 32:35}.
				}
		and he that receiveth light, and continueth in God, receiveth more light;%
			\footnote{%
				% \label{}%
				See \hyperref[DC93-28]{DC 93:28}.
				}
		and that light groweth brighter and brighter until the perfect day.

		\begin{framed}
		\scriptsize
			\begin{compactdesc}
				\item[JSP] that which is of God is light \& he that receiveth light \& continueth in god receiveth more light \& that light groweth \sout{lighter} brighter and brighter untill the perfect day
				\item[Gilbert] that which is of God is \uline{light}, \& he that receiveth light \& continueth in God receiveth more light, and that light groweth brighter \& brighter untill the perfect day.
				\item[Hyde and Smith] That which is of God is light \& he that receivith light \& continueth in God receiveth more light \& that light groweth brighter \& brighter until the perfect day
				\item[BC] 21 That which is of God is light, and he that receiveth light and continueth in God, receiveth more light, and that light groweth brighter and brighter, until the perfect day.
			\end{compactdesc}
		\end{framed}

	% -----
	\subsubsection{Section 50:25} % *DC50-25 
	% -----
		\label{DC50-25}

		And again, verily I say unto you, and I say it that you may know the truth, that you may chase darkness from among you;

		\begin{framed}
		\scriptsize
			\begin{compactdesc}
				\item[JSP] \& again verily I say unto you \& I say it that you may know the truth that you may \sout{choose} chase darkness from among you
				\item[Gilbert] And again verily I say unto you, \& I say it that you may know the truth, that ye may \uline{Chase} \uline{darkness} from among you,
				\item[Hyde and Smith] \& again verily I say unto you \& I say it that you may know the truth that you may chase darkness from among you
				\item[BC] 22 And again, verily I say unto you, and I say it that you may know the truth, that you may chase darkness from among you,
			\end{compactdesc}
		\end{framed}

	% -----
	\subsubsection{Section 50:26} % *DC50-26 
	% -----
		\label{DC50-26}

		He that is ordained of God and sent forth, the same is appointed%
			\footnote{%
				% \label{}%
				See \hyperref[DC29-43]{DC 29:43} and notes there.
				}
		to be the greatest, notwithstanding he is the least and the servant of all.

		\begin{framed}
		\scriptsize
			\begin{compactdesc}
				\item[JSP] for he that is ordaned of God \& sent forth the same is apponted to be the greatest notwithstanding he is least \& the servent of all
				\item[Gilbert] for he that is ordained of God \& sent forth, the same is appointed to be the greatest, notwithstanding he is least \& the servant of all,
				\item[Hyde and Smith] for he that is ordained of God \& sent forth the same is appointed to be greatest nevertheless he is least \& the servant of all
				\item[BC] for he that is ordained of God and sent forth, the same is appointed to be the greatest, notwithstanding he is least, and the servant of all:
			\end{compactdesc}
		\end{framed}

	% -----
	\subsubsection{Section 50:27} % *DC50-27 
	% -----
		\label{DC50-27}

		Wherefore, he is possessor of all things;%
			\footnote{%
				% \label{}%
				Compare \hyperref[DC76-55]{DC 76:55}, and also \hyperref[1CORINTHIANS-15-28]{1 Corinthians 15:28}, \hyperref[HEBREWS-2-8]{Hebrews 2:8}, \hyperref[MORONI-9-26]{Moroni 9:26}, \hyperref[DC63-59]{DC 63:59}, \hyperref[DC121-4]{DC 121:4}, and \hyperref[DC132-20]{DC 132:20}.
				}
		for all things are subject unto him, both in heaven and on the earth, the life and the light, the Spirit and the power, sent forth by the will of the Father through Jesus Christ, his Son.

		\begin{framed}
		\scriptsize
			\begin{compactdesc}
				\item[JSP] wherefore he is possessor of all things <for all things> are subject unto him both in Heaven \& on the Earth the life \& the light the spirit \& the power sent forth by the will of the father through Jesus Christ his son
				\item[Gilbert] wherefore, he is possessor of all things for all things are subject unto him both in Heaven \& on the earth, the life and the light, the \uline{Spirit} \& the \uline{power}, sent forth by the will of the Father through Jesus Christ his son,
				\item[Hyde and Smith] wherefore he is possessor of all things for all things are subject unto him both in heaven \& on the earth the life \& the light the spirit \& the power sent forth by the will of the father through Jesus Christ his son
				\item[BC] 23 Wherefore he is possessor of all things, for all things are subject unto him, both in heaven and on the earth, the life, and the light, the spirit, and the power, sent forth by the will of the Father, through Jesus Christ, his Son;
			\end{compactdesc}
		\end{framed}

	% -----
	\subsubsection{Section 50:28} % *DC50-28 
	% -----
		\label{DC50-28}

		But no man is possessor of all things except he be purified and cleansed from all sin.

		\begin{framed}
		\scriptsize
			\begin{compactdesc}
				\item[JSP] but no man is possessor of all things except he be purified \& cleansed from all sin
				\item[Gilbert] but no man is possessor of all things except he be purified \& Cleansed from all sin,
				\item[Hyde and Smith] but no man is possessor of all things except he be purified and cleansed from all sin
				\item[BC] 24 But no man is possessor of all things, except he be purified and cleansed from all sin;
			\end{compactdesc}
		\end{framed}

	% -----
	\subsubsection{Section 50:29} % *DC50-29 
	% -----
		\label{DC50-29}

		And if ye are purified and cleansed from all sin, ye shall ask whatsoever you will in the name of Jesus and it shall be done.%
			\footnote{%
				% \label{}%
				See versions of this promise \hyperref[ASK-receive]{here}.
				}

		\begin{framed}
		\scriptsize
			\begin{compactdesc}
				\item[JSP] \& if ye are purified \& cleansed from all sin ye shall ask whatsoever you will in the name of Jesus \& it shall be done
				\item[Gilbert] \& if ye are purified \& cleansed from all sin, ye shall ask whatsoever you will in the name of Jesus \& it shall be done,
				\item[Hyde and Smith] \& if ye are purified \& cleansed from all sin ye shall ask whatsoever you will in the name of Jesus \& it shall be done
				\item[BC] 25 And if ye are purified and cleansed from all sin, ye shall ask whatsoever you will in the name of Jesus, and it shall be done:
			\end{compactdesc}
		\end{framed}

	% -----
	\subsubsection{Section 50:30} % *DC50-30 
	% -----
		\label{DC50-30}

		But know this, it shall be given you what you shall ask; and as ye are appointed to the head, the spirits shall be subject unto you.

		\begin{framed}
		\scriptsize
			\begin{compactdesc}
				\item[JSP] but know this it shall be given you what ye shall ask \& as ye are appointed to the head The spirits shall be subject unto you
				\item[Gilbert] but know this it shall be given you what ye shall ask and as ye are appointed to the head the Spirit shall be subject unto you,
				\item[Hyde and Smith] but know this it shall be given you what ye shall ask \& as ye are appointed to the head the spirits shall be subject unto you
				\item[BC] 26 But know this, it shall be given you what you shall ask, and as ye are appointed to the head, the spirits shall be subject unto you:
			\end{compactdesc}
		\end{framed}

	% -----
	\subsubsection{Section 50:31} % *DC50-31 
	% -----
		\label{DC50-31}

		Wherefore, it shall come to pass, that if you behold a spirit manifested that you cannot understand, and you receive not that spirit, ye shall ask of the Father in the name of Jesus; and if he give not unto you that spirit,%
			\footnote{%
				% \label{}%
				Reveal its identity? Tell you what it means?
				}
		then you may know that it is not of God.

		\begin{framed}
		\scriptsize
			\begin{compactdesc}
				\item[JSP] wherefore it shall come to pass that if ye behold a spirit manifested that ye cannot understand \& you receive not that spirit ye shall ask of the father in the name of Jesus \& if he \sout{gave} give not unto you that spirit\sout{s} then ye may know that it is not of God
				\item[Gilbert] wherefore it shall come to pass that if ye behold a spirit manifested that ye cannot understand, and you receive not that Spirit, ye shall ask of the Father in the name of Jesus, \& if he give not unto you that Spirit then you may know that it is not of God,
				\item[Hyde and Smith] wherefore it shall come to pass that if ye behold a spirit manifested that ye cannot understand \& you receive not that spirit ye shall ask of the father in the name of Jesus \& if he give not unto you that spirit then you may know that it is not of God
				\item[BC] 27 Wherefore it shall come to pass, that if you behold a spirit manifested that ye can not understand, and you receive not that spirit, ye shall ask of the Father in the name of Jesus, and if he give not unto you that spirit, then you may know that it is not of God:
			\end{compactdesc}
		\end{framed}

	% -----
	\subsubsection{Section 50:32} % *DC50-32 
	% -----
		\label{DC50-32}

		And it shall be given unto you,%
			\footnote{%
				% \label{}%
				This might be the most inexplicable and unjustifiable comma in all of the DC. None of the older versions have it; its only purpose is to make this passage more difficult to understand. I think they might have included it because the first part of the verse says ``it shall be given'' instead of the expected ``there shall be given,'' but I really don't know. 
				}
		power over that spirit; and you shall proclaim against that spirit with a loud voice that it is not of God---

		\begin{framed}
		\scriptsize
			\begin{compactdesc}
				\item[JSP] \& it shall be given unto you power \sout{that spi} over that spirit \& you shall proclaim against that spirit with a loud voice that it is not of God
				\item[Gilbert] and it shall be given unto you power over that spirit, \& you shall proclaim against that Spirit with a loud voice that \uline{it is not} of \uline{God},
				\item[Hyde and Smith] \& it shall be given unto you power over that spirit \& you shall proclaim against that spirit with a loud voice that it is not of God
				\item[BC] 28 And it shall be given unto you power over that spirit, and you shall proclaim against that spirit with a loud voice, that it is not of God;
			\end{compactdesc}
		\end{framed}

	% -----
	\subsubsection{Section 50:33} % *DC50-33 
	% -----
		\label{DC50-33}

		Not with railing accusation,%
			\footnote{%
				% \label{}%
				See \hyperref[2PETER-2-11]{2 Peter 2:11} and \hyperref[JUDE-1-9]{Jude 1:9}.
				}
		that ye be not overcome, neither with boasting nor rejoicing, lest you be seized therewith.

		\begin{framed}
		\scriptsize
			\begin{compactdesc}
				\item[JSP] not with railing accusation that ye be not over come neither with boasting nor rejoicing lest you be seased [seized] therewith
				\item[Gilbert] not with \uline{railing} \uline{accusation}, that ye be not overcome, neither with \uline{boasting} nor \uline{rejoicing}, lest ye be \uline{seized} therewith,
				\item[Hyde and Smith] \& not with railing accusation that ye be not over come neither with boasting nor rejoiceing lest you be seized therewith
				\item[BC] 29 Not with railing accusation, that ye be not overcome; neither with boasting, nor rejoicing, lest you be seized therewith:
			\end{compactdesc}
		\end{framed}

	% -----
	\subsubsection{Section 50:34} % *DC50-34 
	% -----
		\label{DC50-34}

		He that receiveth of God, let him account it of God; and let him rejoice that he is accounted of God worthy to receive.

		\begin{framed}
		\scriptsize
			\begin{compactdesc}
				\item[JSP] he that receiveth of God let it [him] account it of God \& let him rejoice that he is accounted of God worthy to receive
				\item[Gilbert] He that receiveth of God let him account it of God, \& let him rejoice that he is accounted of God worthy to receive,
				\item[Hyde and Smith] he that receiveth of God let him account it of God \& let him rejoice that he is accounted of God worthy to receive
				\item[BC] 30 He that receiveth of God, let him account it of God, and let him rejoice that he is accounted of God worthy to receive,
			\end{compactdesc}
		\end{framed}

	% -----
	\subsubsection{Section 50:35} % *DC50-35 
	% -----
		\label{DC50-35}
		
			\footnote{%
				% \label{}%
				Tough to parse the connection between this verse and the previous.
				}%
		And by giving heed and doing these things which ye have received, and which ye shall hereafter receive---and the kingdom is given you of the Father, and power to overcome all things which are not ordained of him---

		\begin{framed}
		\scriptsize
			\begin{compactdesc}
				\item[JSP] \& by giving heed \& doing these things which ye have received \& which ye shall hereafter receive \textit{[illegible]} the Kingdom is given unto you of the father \& power to over come all things which is not ordained of him
				\item[Gilbert] \& by giving heed \& doing these things which ye have received \& which ye shall hereafter receive, and the kingdom is given unto you of the Father \& power to overcome all things which is not ordained of him
				\item[Hyde and Smith] \& by giveing heed \& doing these things which ye have received \& which ye shall hereafter receive \& the Kingdom is given unto you of the father and power to overcome all things which is not ordained of him
				\item[BC] and by giving heed and doing these things which ye have received, and which ye shall hereafter receive: 31 And the kingdom is given unto you of the Father, and power to overcome all things, which is not ordained of him:
			\end{compactdesc}
		\end{framed}

	% -----
	\subsubsection{Section 50:36} % *DC50-36 
	% -----
		\label{DC50-36}

		And behold, verily I say unto you, blessed are you who are now hearing these words of mine from the mouth of my servant, for your sins are forgiven you.

		\begin{framed}
		\scriptsize
			\begin{compactdesc}
				\item[JSP] \& Behold verily I say unto you blessed are you that hear these words of mine from the mouth of my servent for your sins are forgiven you
				\item[Gilbert] \& behold verily I say unto you, <blessed are you that> hear these words of mine from the mouth of my Servant for your Sins are forgiven you.---
				\item[Hyde and Smith] \& behold verily I say unto you blessed are you that hear these words of mine from the mouth of my servant for your sins are forgiven you,
				\item[BC] 32 And behold, verily I say unto you, blessed are you who are now hearing these words of mine from the mouth of my servant, for your sins are forgiven you.
			\end{compactdesc}
		\end{framed}

	% -----
	\subsubsection{Section 50:37} % *DC50-37 
	% -----
		\label{DC50-37}

		Let my servant Joseph Wakefield, in whom I am well pleased, and my servant Parley P. Pratt go forth among the churches and strengthen them by the word of exhortation;

		% \begin{framed}
		% \scriptsize
			% \begin{compactdesc}
				% \item[JSP] 
				% % \item[BC] 
				% % \item[]
			% \end{compactdesc}
		% \end{framed}

	% -----
	\subsubsection{Section 50:38} % *DC50-38 
	% -----
		\label{DC50-38}

		And also my servant John Corrill, or as many of my servants as are ordained unto this office, and let them labor in the vineyard; and let no man hinder them doing that which I have appointed unto them---

		% \begin{framed}
		% \scriptsize
			% \begin{compactdesc}
				% \item[JSP] 
				% % \item[BC] 
				% % \item[]
			% \end{compactdesc}
		% \end{framed}

	% -----
	\subsubsection{Section 50:39} % *DC50-39 
	% -----
		\label{DC50-39}

		Wherefore, in this thing my servant Edward Partridge is not justified; nevertheless let him repent and he shall be forgiven.

		% \begin{framed}
		% \scriptsize
			% \begin{compactdesc}
				% \item[JSP] 
				% % \item[BC] 
				% % \item[]
			% \end{compactdesc}
		% \end{framed}

	% -----
	\subsubsection{Section 50:40} % *DC50-40 
	% -----
		\label{DC50-40}

		Behold, ye are little children and ye cannot bear all things now; ye must grow in grace and in the knowledge of the truth.

		% \begin{framed}
		% \scriptsize
			% \begin{compactdesc}
				% \item[JSP] 
				% % \item[BC] 
				% % \item[]
			% \end{compactdesc}
		% \end{framed}

	% -----
	\subsubsection{Section 50:41} % *DC50-41 
	% -----
		\label{DC50-41}

		Fear not, little children, for you are mine, and I have overcome the world, and you are of them that my Father hath given me;

		% \begin{framed}
		% \scriptsize
			% \begin{compactdesc}
				% \item[JSP] 
				% % \item[BC] 
				% % \item[]
			% \end{compactdesc}
		% \end{framed}

	% -----
	\subsubsection{Section 50:42} % *DC50-42 
	% -----
		\label{DC50-42}

		And none of them that my Father hath given me shall be lost.

		% \begin{framed}
		% \scriptsize
			% \begin{compactdesc}
				% \item[JSP] 
				% % \item[BC] 
				% % \item[]
			% \end{compactdesc}
		% \end{framed}

	% -----
	\subsubsection{Section 50:43} % *DC50-43 
	% -----
		\label{DC50-43}

		And the Father and I are one.  I am in the Father and the Father in me; and inasmuch as ye have received me, ye are in me and I in you.

		% \begin{framed}
		% \scriptsize
			% \begin{compactdesc}
				% \item[JSP] 
				% % \item[BC] 
				% % \item[]
			% \end{compactdesc}
		% \end{framed}

	% -----
	\subsubsection{Section 50:44} % *DC50-44 
	% -----
		\label{DC50-44}

		Wherefore, I am in your midst, and I am the good shepherd, and the stone of Israel.  He that buildeth upon this rock shall never fall.

		% \begin{framed}
		% \scriptsize
			% \begin{compactdesc}
				% \item[JSP] 
				% % \item[BC] 
				% % \item[]
			% \end{compactdesc}
		% \end{framed}

	% -----
	\subsubsection{Section 50:45} % *DC50-45 
	% -----
		\label{DC50-45}

		And the day cometh that you shall hear my voice and see me, and know that I am.

		% \begin{framed}
		% \scriptsize
			% \begin{compactdesc}
				% \item[JSP] 
				% % \item[BC] 
				% % \item[]
			% \end{compactdesc}
		% \end{framed}

	% -----
	\subsubsection{Section 50:46} % *DC50-46 
	% -----
		\label{DC50-46}

		Watch, therefore, that ye may be ready.  Even so.  Amen.

		% \begin{framed}
		% \scriptsize
			% \begin{compactdesc}
				% \item[JSP] 
				% % \item[BC] 
				% % \item[]
			% \end{compactdesc}
		% \end{framed}